We investigated whether mechanistic interpretability tools developed for language models transfer to Vision Transformers trained without language supervision, which is a question left largely open by the existing literature. By training Sparse Autoencoders on the MLP activations of a pre-trained ViT-B/16 and using CLIP as an external labelling mechanism, we showed that sparse, interpretable features can be recovered from purely visual representations. \textbf{[TODO --- complete once the final run is available: report here, in one or two sentences, the depth-wise pattern actually observed across the labelled features at the mid-network (Layer 6) and late-network (Layer 11) depths, with concrete example concepts drawn from the CLIP labels and the independent human check (see Results and Analysis) --- e.g.\ a progression from low-level visual properties such as textures and colour patterns to more semantic, object-level concepts. Make sure every named concept belongs to the candidate vocabulary of the dataset actually used for the reported run (see Methodology, \textit{Backbone Model and Dataset}), so that the example is reproducible from the described setup.]} Such a gradient of abstraction, if confirmed, would be consistent with what has been observed in convolutional networks, and our work would characterise it at the level of individual sparse directions in a transformer's residual stream.

The cross-modal evaluation strategy served as a practical solution to the labelling problem that arises when no built-in text-image alignment is available. Rather than depending on a vision-language backbone, CLIP was used as a detached evaluator, decoupling the interpretability analysis from the architecture under study. This choice is replicable for any vision model and does not require re-training or architectural modifications.

Several limitations constrain the scope of these conclusions. The CLIP vocabulary used for labelling is finite and biased towards ImageNet-style concepts, which may cause semantically distinct features to receive the same or similar labels. The analysis is also limited to MLP outputs; the residual stream and attention head outputs represent additional sites of computation that we leave for future work. Finally, the notion of monosemanticity used here is operational rather than formal: as detailed in the Methodology, a feature is labelled monosemantic when visual inspection of its top-$K$ activating crops confirms that they consistently depict the concept identified by CLIP's aggregate label, a qualitative check complemented, for a subset of features, by the independent human assessment reported in Results and Analysis. This remains a small-scale, qualitative proxy that would benefit from a more systematic human validation at scale.

Natural extensions of this work include applying circuit analysis via activation patching to trace how the features identified here contribute causally to specific classification decisions, and comparing the structure of SAE dictionaries learned from a supervised ViT versus a DINO-pretrained ViT---two training regimes that are known to produce qualitatively different internal representations \cite{Raghu2021}.
